\clearpage
\VerdierPhysicalPageAnchor{0264}%
\begingroup
\small
\setlength{\parindent}{0pt}%
\setlength{\parskip}{0pt}%
\setlength{\emergencystretch}{2em}%
\newcommand{\VerdierTermRow}[3]{%
  \VerdierSemanticAnchor{#1}%
  \par\noindent\hangindent=1.5em\hangafter=1
  #2。\enspace\textup{#3。}\par\smallskip}

\VerdierTermRow{idxt-p264-r01}{三角構造と両立する乗法的系}{第 II 章、2.1.2}
\VerdierTermRow{idxt-p264-r02}{射からなる乗法的系}{第 II 章、2.1.1、2.1.4}
\VerdierTermRow{idxt-p264-r03}{飽和乗法的系}{第 II 章、2.1.1}
\VerdierTermRow{idxt-p264-r04}{反完全三角}{第 II 章、1.1.6}
\VerdierTermRow{idxt-p264-r05}{$\mathbb Z$-圏の三角}{第 I 章、3.2.1}
\VerdierTermRow{idxt-p264-r06}{完全三角}{第 II 章、1.1.1}
\VerdierTermRow{idxt-p264-r07}{$K(C)$ の完全三角}{第 I 章、3.3.1}
\VerdierTermRow{idxt-p264-r08}{$\mathbb Z$-圏（それぞれ $\mathbb Z[n]$-圏）}{第 I 章、1.2.2}
\VerdierTermRow{idxt-p264-r09}{厳密な $\mathbb Z$-圏（それぞれ厳密な $\mathbb Z[n]$-圏）}{第 I 章、1.2.2}
\VerdierTermRow{idxt-p264-r10}{$\mathbb Z$-関手、$\mathbb Z[n]$-関手（それぞれ厳密）}{第 I 章、1.2.4}
\VerdierTermRow{idxt-p264-r11}{捩れ $\mathbb Z[n]$-関手}{第 I 章、1.4.4 および第 I 章、1.6.6}
\VerdierTermRow{idxt-p264-r12}{第零コホモロジー対象}{第 II 章、1.3.4}
\VerdierTermRow{idxt-p264-r13}{アーベル圏の対象の $\infty$-$M$-表示}{第 III 章、2.2.1}
\endgroup
